\chapter{Sistemas lineales de EDOS de orden 1}
En este capítulo estudiaremos los sistemas de EDOS de la forma
\[
\begin{cases}
x'_{1}=f_{1}\left(t, x_{1}, \ldots, x_{n}\right) \\
\vdots \\
x_{n}'=f_{n}\left(t, x_{1}, \ldots, x_{n}\right)
\end{cases}
, \; t \in I = \left(\alpha, \omega \right).\]
Si el sistema es lineal, necesariamente debe ser que las funciones $\displaystyle f_{i} $ son lineales, es decir, para $\displaystyle i = 1, \ldots, n $ tendremos que 
\[f_{i}\left(t, x_{1}, \ldots, x_{n}\right)=a_{i1}\left(t\right)x_{1}\left(t\right) + \cdots + a_{in}\left(t\right)x_{n}\left(t\right) .\]
De esta forma, podemos recurrir a la notación matricial
\[\vec{x}'\left(t\right)=A\left(t\right)\vec{x}\left(t\right)+\vec{b}\left(t\right) , \]
donde 
\[\vec{x}\left(t\right)=\begin{pmatrix} x_{1}\left(t\right) \\ \vdots \\ x_{n}\left(t\right) \end{pmatrix}, \quad \vec{x}'\left(t\right)=\frac{d}{dt}\begin{pmatrix} x_{1}\left(t\right) \\ \vdots \\ x_{n}\left(t\right) \end{pmatrix}=\begin{pmatrix} x_{1}'\left(t\right) \\ \vdots \\ x_{n}\left(t\right) \end{pmatrix}  .\]
\[ A\left(t\right)=\begin{pmatrix} a_{11}\left(t\right) & \cdots & a_{1n}\left(t\right) \\ \vdots & & \vdots \\ a_{n1}\left(t\right) & \cdots & a_{nn}\left(t\right) \end{pmatrix}\in \mathcal{M}_{n}\left(t\right) \footnote{Esta notación indica que la matriz depende de $\displaystyle t $.} , \quad \vec{b}\left(t\right)=\begin{pmatrix} b_{1}\left(t\right) \\ \vdots \\ b_{n}\left(t\right) \end{pmatrix}.\]
Decimos que $\displaystyle \vec{b}\left(t\right) $ es el \textbf{vector de términos independientes}. \\
Los objetivos del tema serán:
\begin{itemize}
\item Comprender la estructura del conjunto de soluciones del sistema lineal $\displaystyle \vec{x}'=A\vec{x}+\vec{b} $ y el del sistema lineal homogéneo asociado $\displaystyle \vec{x}' = A\vec{x} $. Para este resultado es fundamental la siguiente versión del teorema de existencia y unicidad del problema de Cauchy.
\item Aprender a resolver los sistemas $\displaystyle \vec{x}'=A\vec{x} + \vec{b} $, donde $\displaystyle A \in \mathcal{n}\left(\R\right) $, es decir, tiene coeficientes constantes.
\item Veremos el caso particular de las ecuaciones de orden $\displaystyle n $, que son de la forma
	\[x^{\left(n\right)} + a_{n-1}x^{\left(n-1\right)} + \cdots + a_{1}x' + a_{0}x + b = 0 ,\]
	donde $\displaystyle A $ tiene un 'aspecto' determinado \footnote{En el primer tema vimos que las EDOS de orden $\displaystyle n $ se puede escribir como un sistema de orden 1.}.
\item Hacer análisis cualitativo de los sistemas autónomos, es decir, $\displaystyle \vec{x}' = A\vec{x} $ con $\displaystyle A $ de coeficientes constantes y estamos en el caso de dimensión 2, es decir, $\displaystyle \vec{x}=\left(x_{1}, x_{2}\right)^{t} $.
\end{itemize}
\begin{notation}
Utilizaremos la siguiente notación
\[\int^{t}_{t_{0}} \vec{x}\left(s\right) \; d s = \begin{pmatrix} \int^{t}_{t_{0}} x_{1}\left(s\right) \; d s \\ \vdots \\ \int^{t}_{t_{0}} x_{n}\left(s\right) \; d s \end{pmatrix} .\]
Recordamos que $\displaystyle A\left(t\right) $ y $\displaystyle \vec{b}\left(t\right) $ son continuas si y solo si sus componentes lo son.
\end{notation}

\begin{theorem}[Teorema de existencia y unicidad]
Si $\displaystyle A\left(t\right) $ y $\displaystyle \vec{b}\left(t\right) $ son continuas en un intervalo $\displaystyle I = \left(\alpha, \omega \right)\subset \R $, entonces el problema de Cauchy
\[
\begin{cases}
\vec{x}' = A\left(t\right)\vec{x} + \vec{b}\left(t\right)\\
\vec{x}\left(t_{0}\right) = \vec{x}_{0}, \; t_{0} \in \left(\alpha, \omega\right)
\end{cases}
\]
tiene solución única $\displaystyle \vec{x}\left(t\right) $ en $\displaystyle \left(\alpha, \omega\right) $ . 
\end{theorem}
\section{Estructura de las soluciones}
Asumimos que $\displaystyle A\left(t\right) $ y $\displaystyle \vec{b}\left(t\right) $ son continuas en $\displaystyle I = \left(\alpha, \omega \right) \subset \R $. 
\begin{theorem}
El conjunto de soluciones del sistema $\displaystyle \vec{x}'=A\left(t\right)\vec{x} $ tiene estructura  de $\displaystyle \R $-espacio vectorial de dimensión $\displaystyle n $. 
\end{theorem}
\begin{proof}
El conjunto de soluciones tendrá estructura de espacio vectorial puesto que si $\displaystyle \vec{x}_{1}\left(t\right) $ y $\displaystyle \vec{x}_{2}\left(t\right) $ son soluciones y $\displaystyle \alpha, \beta \in \R $ se tiene que 
\[\left(\alpha\vec{x}_{1} + \beta\vec{x}_{2}\right)' = \alpha\vec{x}_{1}' + \beta\vec{x}_{2}' = \alpha A\left(t\right)\vec{x}_{1}+\beta A\left(t\right)\vec{x}_{2}= A\left(t\right)\left(\alpha \vec{x}_{1} + \vec{x}_{2}\right) .\]
De esta forma, $\displaystyle \alpha\vec{x}_{1} + \beta\vec{x}_{2} $ también es solución del sistema. Veamos ahora que tiene dimensión $\displaystyle n $. Consideremos los siguientes $\displaystyle n $ problemas de valor inicial:
\[\left(P_{i}\right)
\begin{cases}
\vec{x}'=A\left(t\right)\vec{x} \\
\vec{x}\left(t_{0}\right)=\vec{e}_{i}
\end{cases}
, \; t_{0} \in I, \; i = 1, \ldots, n.\]
Por el teorema de existencia y unicidad, cada $\displaystyle P_{i} $ tiene solución única $\displaystyle \vec{x}_{i}\left(t\right) $. Consideramos el conjunto $\displaystyle \left\{ \vec{x}_{1}, \ldots, \vec{x}_{n}\right\}  $ y veamos que es una base del conjunto de soluciones del sistema. Veamos primero que forma un sistema de generadores. Sea $\displaystyle \vec{z}\left(t\right) $ una solución cualquiera del sistema, en particular $\displaystyle \vec{z}\left(t_{0}\right)=\vec{z}_{0} $. Como $\displaystyle \left\{ \vec{e}_{1}, \ldots, \vec{e}_{n}\right\}  $ es base de $\displaystyle \R^{n} $ tendremos que existen $\displaystyle c_{1}, \ldots, c_{n} \in \R $ tales que
\[\vec{z}_{0} =  c_{1}\vec{e}_{1} + \cdots + c_{n}\vec{e}_{n}.\]
Entonces tenemos que $\displaystyle c_{1}\vec{x}_{1} + \cdots + c_{n}\vec{x}_{n} $ es solución del problema de Cauchy
\[
\begin{cases}
\vec{x}' = A\left(t\right)\vec{x} \\
\vec{x}\left(t_{0}\right)=z_{0}
\end{cases}
.\]
Por el teorema de existencia y unicidad debe ser que $\displaystyle \vec{z} = c_{1}\vec{x}_{1} + \cdots + c_{n}\vec{x}_{n}$. Ahora, veamos que son linealmente independientes. Supongamos que existen $\displaystyle a_{1}, \ldots, a_{n} \in \R $ tales que 
\[a_{1}\vec{x}_{1} + \cdots + a_{n}\vec{x}_{n}=0, \; \forall t \in I = \left(\alpha , \omega \right) .\]
En particular, para $\displaystyle t_{0}  $ tendremos que
\[a_{1}\vec{x}_{1}\left(t_{0}\right) +\cdots + a_{n}\vec{x}_{n}\left(t_{0}\right)=0 \Rightarrow a_{1}\vec{e}_{1} + \cdots + a_{n}\vec{e}_{n}= 0 \Rightarrow a_{1} = \cdots =a_{n}=0 .\]
Así, tenemos que $\displaystyle \left\{ \vec{x}_{1}, \ldots, \vec{x}_{n}\right\}  $ es una base del conjunto de soluciones del sistema. 
\end{proof}

\begin{colorary}
La solución general del sistema completo $\displaystyle \vec{x}' = A\left(t\right)\vec{x} +\vec{b}\left(t\right)$ es 
\[\vec{x}\left(t\right)=\vec{x}_{h}\left(t\right)+\vec{x}_{p}\left(t\right) ,\]
donde $\displaystyle \vec{x}_{h} $ es la solución general del sistema homogéneo asociado $\displaystyle \vec{x}' = A\left(t\right)\vec{x} $, y $\displaystyle \vec{x}_{p} $ es una solución particular del completo. 
\end{colorary}
\begin{proof}
Sea $\displaystyle \vec{x}_{h} $ una solución del sistema homogéneo asociado y sea $\displaystyle \vec{x}_{p} $ una solución particular del sistema completo. Entonces,
\[\left(\vec{x}_{h}+\vec{x}_{p}\right)' = \vec{x}_{h}' + \vec{x}_{p}' = A\left(t\right)\vec{x}_{h}+A\left(t\right)\vec{x}_{p} + \vec{b}\left(t\right)=A\left(t\right)\left(\vec{x}_{h}+\vec{x}_{p}\right)+\vec{b}\left(t\right)	 .\]
Así, tenemos que $\displaystyle \vec{x}_{h}+\vec{x}_{p} $ es solución del sistema completo. Ahora, veamos que toda solución del sistema se puede escribir de esta forma. Dada una solución $\displaystyle \vec{x} $ del sistema y una solución particular $\displaystyle \vec{x}_{p} $, tenemos que
\[\left(\vec{x}-\vec{x}_{p}\right)' = \vec{x}' - \vec{x}_{p}' = A\left(t\right)\vec{x} + \vec{b}\left(t\right) - \left(A\left(t\right)\vec{x}_{p}+\vec{b}\left(t\right)\right)=A\left(t\right)\left(\vec{x}-\vec{x}_{p}\right).\]
Así, $\displaystyle \vec{x}-\vec{x}_{p} $ es solución del sistema homogéneo asociado. 
\end{proof}
\section{Matriz fundamental}

\begin{definition}[Matriz fundamental]
Dado un sistema lineal homogéneo $\displaystyle \vec{x}' = A\left(t\right)\vec{x} $ con $\displaystyle A \in C\left(I, \mathcal{M}_{n}\right) $ \footnote{Esta notación indica que $\displaystyle A $ es continua en el intervalo $\displaystyle I $ y es una matriz cuadrada de dimensión $\displaystyle n $.}, denominamos \textbf{matriz fundamental} a aquella cuyas columnas son $\displaystyle n $ soluciones linealmente independientes del sistema y la denotamos $\displaystyle \phi\left(t\right) $. 
\end{definition}
\begin{observation}
Por el teorema anterior tendremos que $\displaystyle \phi\left(t\right) $ es también una matriz de dimensión $\displaystyle n \times n $ y sus entradas son funciones continuas, por lo que $\displaystyle \phi\left(t\right)\in C\left(I, \mathcal{M}_{n}\right) $. 
\end{observation}
\begin{notation}
	A cada columna de $\displaystyle \phi\left(t\right) $ la denotamos $\displaystyle \phi^{i}\left(t\right) $ para $\displaystyle i= 1, \ldots, n $. Así, tendremos que $\displaystyle \left\{ \phi^{1}\left(t\right), \ldots, \phi^{n}\left(t\right)\right\}  $ es una base de las soluciones $\displaystyle \vec{x}' = A\left(t\right)\vec{x} $. 
\end{notation}
\begin{eg}
	Si $\displaystyle t_{0} \in I $ y $\displaystyle \vec{x}_{i} $ es solución del sistema $\displaystyle \vec{x}' = A\left(t\right)\vec{x} $ tal que $\displaystyle \vec{x}_{i}\left(t_{0}\right)=e_{i} $ para $\displaystyle i \in \left\{ 1, \ldots, n\right\}  $, entonces $\displaystyle \phi\left(t\right)=\left(\vec{x}_{1}, \ldots, \vec{x}_{n}\right) $ es una matriz fundamental del sistema \footnote{Esto se deduce de la demostración del teorema anterior.}. 
\end{eg}

\begin{prop}
Supongamos que $\displaystyle \phi\left(t\right) $ es matriz fundamental del sistema $\displaystyle \vec{x}' = A\left(t\right)\vec{x} $ e $\displaystyle \vec{y}\left(t\right) $ es solución del sistema. Entonces, existe $\displaystyle \vec{c}= \left(c_{1}, \ldots, c_{n}\right)\in \R^{n} $ tal que $\displaystyle \vec{y}\left(t\right)=\phi\left(t\right)\vec{c} $. 
\end{prop}
\begin{proof}
Como $\displaystyle \phi\left(t\right) $ es matriz fundamental, sus columnas son linealmente independientes y forman una base de las soluciones del sistema. Entonces, existirán escalares tales que 
\[\vec{y}\left(t\right)=c_{1}\phi^{1}\left(t\right) + \cdots + c_{n}\phi^{n}\left(t\right) .\]
\end{proof}
\begin{prop}
Dado el sistema $\displaystyle \vec{x}' = A\left(t\right)\vec{x} $, $\displaystyle t \in I $ con $\displaystyle A\left(t\right)\in C\left(I, \mathcal{M}_{n}\right) $. Sea $\displaystyle \phi\left(t\right)\in C\left(I, \mathcal{M}_{n}\right) $ tal que sus columnas son soluciones de dicho sistema, entonces son equivalentes:
\begin{enumerate}
\item Existe $\displaystyle t_{0} \in I $ tal que $\displaystyle \det\left(\phi\left(t_{0}\right)\right)\neq 0 $. 
\item $\displaystyle \forall t \in I $, $\displaystyle \det\left(\phi\left(t\right)\right) \neq 0 $. 
\item $\displaystyle \phi\left(t\right) $ es matriz fundamental del sistema. 
\end{enumerate}
\end{prop}
\begin{proof}
\begin{description}
	\item[(1) $\displaystyle \Rightarrow $ (3)] Supongamos que existe $\displaystyle t_{0} \in I $ tal que $\displaystyle \det\left(\phi\left(t_{0}\right)\right)\neq 0 $, entonces los vectores $\displaystyle \left\{ \phi^{1}\left(t_{0}\right), \ldots , \phi^{n}\left(t_{0}\right)\right\}  $ son linealmente independientes. Queremos ver que $\displaystyle \phi\left(t\right) $ es matriz fundamental. Sus columnas son soluciones del sistema $\displaystyle \vec{x}' = A\left(t\right)\vec{x} $ por hipótesis. Quiero ver que $\displaystyle \left\{ \phi^{1}\left(t\right), \ldots, \phi^{n}\left(t\right)\right\}  $ son linealmente independientes. 
		Si existen $\displaystyle c_{1}, \ldots, c_{n} \in \R $ tales que 
		\[c_{1}\phi^{1}\left(t\right) + \cdots + c_{n}\phi^{n}\left(t\right)=0 \Rightarrow c_{1}\phi^{1}\left(t_{0}\right)+\cdots + c_{n}\phi^{n}\left(t_{0}\right)=0 \Rightarrow c_{1} = \cdots = c_{n}= 0 .\]	
\item[(3) $\displaystyle \Rightarrow $ (2)] Supongamos que $\displaystyle \phi\left(t\right) $ es matriz fundamental y que existe $\displaystyle t_{0} \in I $ tal que $\displaystyle \det\left(\phi\left(t_{0}\right)\right)=0 $. 
	Esto quiere decir que $\displaystyle \left\{ \phi^{1}\left(t_{0}\right), \ldots, \phi^{n}\left(t_{0}\right)\right\}  $ son linealmente dependientes, por lo que existen $\displaystyle c_{1}, \ldots, c_{n} \in \R $ no todos nulos tales que 
	\[c_{1}\phi_{1}\left(t_{0}\right)+\cdots + c_{n}\phi_{n}\left(t_{0}\right)=0 .\]
Consideramos $\displaystyle \vec{x}\left(t\right)=c_{1}\phi^{1}\left(t\right) + \cdots + c_{n}\phi^{n}\left(t\right) $. Como $\displaystyle \phi^{i}\left(t\right) $ es solución del sistema para cada $\displaystyle i = 1, \ldots, n $, entonces $\displaystyle \vec{x}\left(t\right) $ también es solución del sistema y además verifica la condición inicial $\displaystyle \vec{x}\left(t_{0}\right)=0 $. 
El problema de valor inicial
\[
\begin{cases}
\vec{x}' = A\left(t\right)\vec{x} \\
\vec{x}\left(t_{0}\right)=0
\end{cases}
,\]
tiene como solución la trivial. Por unicidad de soluciones tendremos que 
\[c_{1}\phi^{1}\left(t\right) + \cdots +c_{n}\phi^{n}\left(t\right)=0 ,\]
con $\displaystyle c_{i} $ no todos nulos. De esta forma, los $\displaystyle \phi^{i}\left(t\right) $ son linealmente dependientes, que es una contradicción. 
\item[(2) $\displaystyle \Rightarrow $ (1)] Trivial. 
\end{description}
\end{proof}
\begin{observation}
Con este resultado se garantiza tener solución en el intervalo entero, es decir, no pueden existir puntos dentro de ese intervalo en el que las soluciones dejen de ser linealmente independientes cuando sé que ya lo son en un cierto punto del intervalo.
\end{observation}

\begin{notation}
De ahora en adelante, denotamos por $\displaystyle \left(S_{h}\right) $ a la EDO homogénea $\displaystyle \vec{x}'=A\left(t\right)\vec{x} $.
\end{notation}
 
\begin{colorary}
Supongamos que $\displaystyle \phi\left(t\right) $ es matriz fundamental del sistema $\displaystyle \vec{x}'=A\left(t\right)\vec{x} $ con $\displaystyle A\left(t\right) \in C\left(I, \mathcal{M}_{n}\right) $. Entonces, la solución del problema de valor inicial
\[ \left(P\right)
\begin{cases}
\vec{x}'=A\left(t\right)\vec{x} \\
\vec{x}\left(t_{0}\right)=\vec{x}_{0}
\end{cases}
, \; t_{0} \in I,\]
es $\displaystyle \vec{x}\left(t\right)=\phi\left(t\right)\phi^{-1}\left(t_{0}\right)\vec{x}_{0} $ 
\end{colorary}
\begin{proof}
Por lo visto anteriormente, tendremos que $\displaystyle \forall t \in I $, $\displaystyle \phi\left(t\right) $ es invertible. También vimos que las soluciones del problema son de la forma $\displaystyle \phi\left(t\right)\vec{c} $ para $\displaystyle \vec{c} \in \R^{n} $. Como queremos que $\displaystyle \vec{x}\left(t\right) $ satisfazca la condición inicial,
\[\vec{x}\left(t_{0}\right)=\vec{x}_{0}=\phi\left(t_{0}\right)\vec{c} \iff \vec{c}=\phi^{-1}\left(t_{0}\right)\vec{x}_{0} .\]
Por tanto, la única solución de $\displaystyle \left(P\right) $ es
\[\vec{x}\left(t\right)=\phi\left(t\right)\phi^{-1}\left(t_{0}\right)\vec{x}_{0} .\]
\end{proof}
\begin{observation}
Por tanto, para resolver explícitamente el sistema $\displaystyle \vec{x}' = A\left(t\right)\vec{x} $ basta con hallar la matriz fundamental. 
\end{observation}

\begin{prop}[Cambio de base]
	Supongamos que $\displaystyle \phi\left(t\right) $ es matriz fundamental de $\displaystyle \left(S_{h}\right) $ con $\displaystyle A\left(t\right)\in C\left(I, \mathcal{M}_{n}\right) $.
	\begin{enumerate}
	\item Sea $\displaystyle C \in \mathcal{M}_{n}\left(\R\right) $, de coeficientes constantes, con $\displaystyle \det\left(C\right)\neq0 $. Entonces, $\displaystyle \phi\left(t\right)C $ es también matriz fundamental.
	\item Si $\displaystyle \psi\left(t\right) $ es otra matriz fundamental del sistema, entonces existe $\displaystyle C \in \mathcal{M}_{n}\left(\R\right) $ con $\displaystyle \det\left(C\right)\neq 0 $ tal que 
	\[\psi\left(t\right)=\phi\left(t\right)C .\]
	\end{enumerate}
\end{prop}
\begin{proof} Supongamos que $\displaystyle \phi\left(t\right) $ es matriz fundamental de $\displaystyle \left(S_{h}\right) $ con $\displaystyle A\left(t\right)\in C\left(I, \mathcal{M}_{n}\right) $. 
\begin{enumerate}
\item Observamos que 
\[\phi\left(t\right)C = \left(\phi\left(t\right)C_{1}, \ldots, \phi\left(t\right)C_{n}\right) ,\]
donde $\displaystyle C_{i} $ son las columnas de $\displaystyle C $. Es fácil ver que $\displaystyle \phi\left(t\right)C_{i} $ es solución de $\displaystyle \left(S_{h}\right) $, $\displaystyle \forall i = 1, \ldots, n $, por ser combinación lineal de soluciones de $\displaystyle \left(S_{h}\right) $.
Ahora veamos que son linealmente independientes. Para $\displaystyle t \in I $, tenemos que 
\[\det\left(\phi\left(t\right)C\right)=\det\left(\phi\left(t\right)\right)\det\left(C\right)\neq 0 .\]
Por tanto, tenemos que $\displaystyle \phi\left(t\right)C $ es matriz fundamental.
\item Para la segunda parte, como $\displaystyle \left\{ \phi^{1}\left(t\right), \ldots, \phi^{n}\left(t\right)\right\}  $ es una base del conjunto de soluciones de $\displaystyle \left(S_{h}\right) $, cada columna de $\displaystyle \psi\left(t\right) $ se puede escribir de la forma
\[\psi^{i}\left(t\right)=c_{1i}\phi^{1}\left(t\right)+ \cdots + c_{ni}\phi^{n}\left(t\right) ,\]
Así, tenemos que $\displaystyle \psi_{i}\left(t\right)=\phi\left(t\right)\vec{c}_{i} $ con $\displaystyle \vec{c}_{i}\in \R^{n} $, que es equivalente a que $\displaystyle \psi\left(t\right)=\phi\left(t\right)C $ con $\displaystyle \det\left(C\right)\neq 0 $. En efecto, si tuviésemos $\displaystyle \det\left(C\right)=0 $, $\displaystyle \forall t \in I $ se tendría que
\[\det\left(\psi\left(t\right)\right)=\det\left(\phi\left(t\right)C\right)=\det\left(\phi\left(t\right)\right)\det\left(C\right)=0,\]
por lo que $\displaystyle \psi\left(t\right) $ no sería matriz fundamental, que es una contradicción. 
\end{enumerate}
\end{proof}
\begin{definition}[Matriz fundamental principal]
Dada una matriz fundamental asociada a $\displaystyle \left(S_{h}\right) $ con $\displaystyle A\left(t\right)\in C\left(I, \mathcal{M}_{n}\right) $, diremos que es \textbf{matriz fundamental principal} o \textbf{canónica} asociada al instante inicial $\displaystyle t_{0} \in I $, si $\displaystyle \psi\left(t_{0}\right)=I$.
\end{definition}
\begin{observation}
Podemos observar que por el teorema de existencia y unicidad podemos garantizar que la matriz fundamental principal existe y es única. \\
Por otro lado, la matriz fundamental principal es útil puesto que simplifica los cálculos para resolver el problema
\[
\begin{cases}
\vec{x}' = A\left(t\right)\vec{x} + \vec{b}\left(t\right) \\
\vec{x}\left(t_{0}\right)=\vec{x}_{0}
\end{cases}
.\]
En efecto, por un corolario anterior tendremos que la solución será  
\[ \vec{x}\left(t\right)=\varphi\left(t\right)\varphi^{-1}\left(t_{0}\right)\vec{x}_{0}=\varphi\left(t\right)\vec{x}_{0}.\]
\end{observation}
\begin{prop}
Sea $\displaystyle \varphi\left(t\right) $ matriz fundamental principal para el problema
\[
\begin{cases}
\vec{x}' = A\left(t\right)\vec{x} \\
\vec{x}\left(t_{0}\right)=\vec{x}_{0}
\end{cases}
,\]
con $\displaystyle A\left(t\right)=A $, es decir de coeficientes constantes, y $\displaystyle t_{0}=0 $. Entonces
\begin{enumerate}
\item $\displaystyle \varphi\left(t+s\right)=\varphi\left(t\right)\varphi\left(s\right) $, $\displaystyle \forall t,s \in \R $. 
\item $\displaystyle \varphi^{-1}\left(t\right)=\varphi\left(-t\right) $, $\displaystyle \forall t \in \R $. 
\end{enumerate}
\end{prop}
\section{Método de variación de constantes}
\begin{theorem}
El problema de valor inicial
\[
\begin{cases}
\vec{x}' = A\left(t\right)\vec{x} + \vec{b}\left(t\right) \\
\vec{x}\left(t_{0}\right)=\vec{x}_{0}, \; t_{0} \in I
\end{cases}
,\]
con $\displaystyle A\left(t\right) $ y $\displaystyle \vec{b}\left(t\right) $ continuas en $\displaystyle I = \left(\alpha, \omega \right) $, tiene por solución única
\[\vec{x}\left(t\right)=\phi\left(t\right)\phi^{-1}\left(t_{0}\right)\vec{x}_{0} + \phi\left(t\right)\int^{t}_{t_{0}} \phi^{-1}\left(s\right)\vec{b}\left(s\right) \; d s ,\]
donde $\displaystyle \phi\left(t\right) $ es matriz fundamental del sistema homogéneo asociado. 
\end{theorem}
\begin{proof}
Dado un sistema lineal no homogéneo $\displaystyle \vec{x}' = A\left(t\right)\vec{x} + \vec{b}\left(t\right) $ y una matriz fundamental $\displaystyle \phi\left(t\right) $ de $\displaystyle \left(S_{h}\right) $, sabemos que $\displaystyle \forall \vec{c} \in \R^{n} $ el vector $\displaystyle \phi\left(t\right)\vec{c} $ es solución de $\displaystyle \left(S_{h}\right) $. \\
Ahora, podemos generalizar el método de variación de constantes que vimos en el tema anterior para encontrar una solución particular del sistema no homogéneo. Concretamente, buscamos una solución de la forma
\[\vec{x}_{p}\left(t\right)=\phi\left(t\right)\vec{c}\left(t\right) .\]
Si forzamos para que $\displaystyle \vec{x}_{p} $ sea solución:
\[\phi'\left(t\right)\vec{c}\left(t\right) + \phi\left(t\right)\vec{c}'\left(t\right)=A\left(t\right)\phi\left(t\right)\vec{c}\left(t\right)+b\left(t\right)=\phi'\left(t\right)\vec{c}\left(t\right)+\vec{b}\left(t\right)  .\]
De aquí obtenemos que $\displaystyle \phi\left(t\right)\vec{c}'\left(t\right)=\vec{b}\left(t\right) $, por lo que $\displaystyle \vec{c}'\left(t\right)=\phi^{-1}\left(t\right)\vec{b}\left(t\right) $. Así, tendremos que
\[\vec{c}\left(t\right)=\int^{t}_{t_{0}} \phi^{-1}\left(s\right)\vec{b}\left(s\right) \; d s .\]
Así, obtenemos la solución particular será
\[\vec{x}_{p}\left(t\right)=\phi\left(t\right)\int^{t}_{t_{0}} \phi^{-1}\left(s\right)\vec{b}\left(s\right) \; d s .\]
La fórmula del teorema se deduce de un teorema anterior que indica la estructura de las soluciones para este tipo de EDOS. 
\end{proof}
\begin{observation}
Podemos ver que no importa la matriz fundamental que escojamos.
\end{observation}

\section{Sistemas lineales con coeficientes constantes}

A partir de ahora consideramos $\displaystyle A\left(t\right)=A $ matriz de coeficientes constantes.
El problema $\displaystyle \vec{x}'=A\vec{x} $ es muy similar al problema $\displaystyle x' = ax $, que tiene como solución general $\displaystyle x\left(t\right)=Ce^{at} $, $\displaystyle C \in \R $. Veamos que el sistema tiene una solución parecida, concretamente, veremos que $\displaystyle \phi\left(t\right)=e^{At} $. 
\subsection{$A$ es diagonalizable}

En primer lugar nos podemos preguntar por las condiciones que deben verificar $\displaystyle \vec{v}\in \R^{n} $ y $\displaystyle \lambda \in \R $ para que $\displaystyle e^{\lambda t}\vec{v} $ sea solución del sistema $\displaystyle \left(S_{h}\right)\; \vec{x}'=A\vec{x} $. Tendremos que 
\[\vec{x}'=A\vec{x} \iff \lambda e^{\lambda t}\vec{v}=Ae^{\lambda t}\vec{v} \iff A\vec{v}-\lambda\vec{v}=0 \iff \left(A-\lambda I\right)\vec{v}=0 .\]
Así, tenemos que $\displaystyle e^{\lambda t}\vec{v} $ es solución si y solo si $\displaystyle \lambda \in \R $ es autovalor de $\displaystyle A $ y $\displaystyle \vec{v} $ es un vector asociado a este autovalor. Por tanto, si $\displaystyle A $ es diagonalizable ya sabemos encontrar una matriz fundamental.
\begin{prop}
	Si $\displaystyle \left\{ \vec{v}_{1}, \ldots, \vec{v}_{n}\right\}  $ son $\displaystyle n $ autovectores linealmente independientes asociados a $\displaystyle \lambda_{1}, \ldots, \lambda_{n} $ autovalores de la matriz $\displaystyle A $, y definimos las matrices
	\[P = \left(\vec{v}_{1}, \ldots, \vec{v}_{n}\right) \quad \text{y} \quad D\left(t\right)=\begin{pmatrix} e^{\lambda_{1}t} & & & \\
	& e^{\lambda_{2}t} & & \\
	& & & \\
	& & & e^{\lambda_{n}t}\end{pmatrix} ,\]
	entonces $\displaystyle \phi\left(t\right)=PD\left(t\right) $ es matriz fundamental de $\displaystyle \left(S_{h}\right) $. Además, la matriz fundamental principal asociada al instante $\displaystyle t_{0}=0 $ es $\displaystyle \psi\left(t\right)=PD\left(t\right)P^{-1} $.
\end{prop}
\begin{proof}
	Tenemos que $\displaystyle \phi\left(t\right)=PD\left(t\right) $ es matriz fundamental por la reflexión que hicimos anteriormente y puesto que por hipótesis $\displaystyle \left\{ \vec{v}_{1}, \ldots, \vec{v}_{n}\right\}  $ son linealmente independientes. Así, las columnas de $\displaystyle \phi $ son $\displaystyle \left\{ e^{\lambda_{1}t}\vec{v}_{1}, \ldots, e^{\lambda_{n}}\vec{v}_{n}\right\}  $, que son soluciones linealmente independientes. 
Por otro lado, 
\[\psi\left(0\right)=PD\left(0\right)P^{-1}=P P^{-1}= I .\]
\end{proof}
\begin{observation}
Si el instante inicial es $\displaystyle t_{0}\neq 0 $, tendremos que la matriz principal será
\[\psi\left(t\right)=P\begin{pmatrix} e^{\lambda_{1}\left(t-t_{0}\right)} & & & \\
	& e^{\lambda_{2}\left(t-t_{0}\right)} & & \\
	& & & \\
	& & & e^{\lambda_{n}\left(t-t_{0}\right)}\end{pmatrix} .\]
\end{observation}
Ahora nos preguntamos, qué pasa si $\displaystyle A $ no es diagonalizable. Así, retomamos la primera idea de definir $\displaystyle \phi\left(t\right) $ de la forma $\displaystyle e^{At} $. 
\subsection{Exponencial de una matriz}
La forma en la que definimos el exponencial de una matriz es extendiendo la definición de serie de potencias de la función exponencial
\[e^{x}=1 + x + \frac{x^{2}}{2!} + \cdots = \sum^{\infty}_{k = 0}\frac{x^{k}}{k!} .\]
\begin{definition}[Exponencial de una matriz]
Dada $\displaystyle A \in \mathcal{M}_{n} $, se define el \textbf{exponencial de $\displaystyle A $} como
\[\exp\left(A\right)=e^{A}:= I + A + \frac{A^{2}}{2!} + \cdots = \sum^{\infty}_{k = 0}\frac{A^{k}}{k!} .\]
\end{definition}
\begin{theorem}
Dado el sistema $\displaystyle \vec{x}' = A \vec{x} $, la matriz fundamental principal tal que $\displaystyle \psi\left(0\right)=I $ viene dada por
\[\psi\left(t\right)=\exp\left(At\right)=\sum^{\infty}_{k = 0}\frac{t ^{k}A^{k}}{k!}, \; t \in \R .\]
Además, la convergencia de la serie es uniforme sobre cada intervalo cerrado y acotado $\displaystyle \left[a,b\right] \subset \R $. 
\end{theorem}
\begin{observation}
Los dos siguientes problemas son equivalentes
\[
\begin{cases}
\phi'\left(t\right)=A\phi\left(t\right) \\
\phi\left(0\right)=I
\end{cases}
\]
y los $\displaystyle n $ problemas de valor inicial de la forma 
\[
\begin{cases}
\vec{x}' = A\vec{x} \\
\vec{x}\left(0\right)=e_{i}, \; i = 1, \ldots, n
\end{cases}
.\]
De esta forma, está claro que el PVI también tiene una única solución $\displaystyle \phi\left(t\right) $.
\end{observation}

\begin{proof}
Si aplicamos el método de las iteradas de Picard tenemos que
\[\psi_{0}\left(t\right)=I .\]
\[\psi_{1}\left(t\right)=I + \int^{t}_{0} AI \; d s = I + tA .\]
\[\psi_{2}\left(t\right)=I + \int^{t}_{0} A\left(I + sA\right) \; d s= I + tA + \frac{t ^{2}}{2}A^{2} .\]
\[\vdots .\]
\[\psi_{j}\left(t\right)=I + tA + \cdots + \frac{t ^{j}}{j!}A^{j} .\]
En un ejercicio del tema 1 vimos que este método converge, por lo que sabemos que existe una matriz $\displaystyle \phi\left(t\right) $ tal que 
\[\lim_{j  \to \infty}\psi_{j}\left(t\right)=\phi\left(t\right), \; t \in \R .\]
Además esta convergencia es uniforme en cada intervalo $\displaystyle \left[a,b\right]  $. Por otro lado, la demostración del teorema de existencia y unicidad (la construcción de las iterantes de Picard) nos asegura que las columnas de $\displaystyle \phi\left(t\right) $ son solución del sistema $\displaystyle \vec{x}' = A\vec{x} $.  
\end{proof}
\begin{observation}
Es obvio que $\displaystyle \phi\left(0\right)=I $. De esta forma, la solución del PVI
\[
\begin{cases}
\vec{x}' = A\vec{x}\\
\vec{x}\left(0\right)=\vec{x}_{0}
\end{cases}
\]
es $\displaystyle \vec{x}\left(t\right)=e^{tA}\vec{x}_{0} $.
\end{observation}
\begin{observation}
Este resultado es coherente con la proposición que acabamos de ver. Veamos que si $\displaystyle A $ es diagonalizable, $\displaystyle \psi\left(t\right)=PD\left(t\right)P^{-1} $ coincide $\displaystyle \exp\left(At\right)$. Supongamos que $\displaystyle A = PDP^{-1} $ con 
	\[P = \left\{ \vec{v}_{1}, \ldots, \vec{v}_{n}\right\} \quad \text{y} \quad D = \begin{pmatrix} \lambda_{1} & &0 & \\
	& \lambda_{2} & & \\
	& 0& & \\
	& & & \lambda_{n}\end{pmatrix} .\]
Entonces, 
\[
\begin{split}
	\exp\left(At\right) = & I + tA + \frac{t ^{2}}{2!}A^{2} + \cdots + \frac{t ^{j}}{j!}A^{j} + \cdots \\
	= &  P P^{-1} + t P D P^{-1} + \frac{t ^{2}}{2!}PD^{2}P^{-1} + \cdots + \frac{t ^{j}}{j!}PD^{j}P^{-1} + \cdots\\
 = & P\begin{pmatrix} \sum^{\infty}_{k = 0}\frac{t ^{k}}{k!}\lambda_{1}^{k} & & & \\
	& \sum^{\infty}_{k = 0}\frac{t ^{k}}{k!}\lambda_{2}^{k} & & \\
	& & & \\
	& & & \sum^{\infty}_{k = 0}\frac{t ^{k}}{k!}\lambda_{n}^{k}\end{pmatrix} P^{-1}\\
	= & P\begin{pmatrix} e^{\lambda_{1}t} & & & \\
	& e^{\lambda_{2}t} & & \\
	& & & \\
	& & & e^{\lambda_{n}t}\end{pmatrix}P^{-1}=PD\left(t\right)P^{-1} = \psi\left(t\right).
\end{split}
\]
\end{observation}
\begin{observation}
Por otro lado, si quisiéramos calcular la matriz fundamental principal asociada al instante $\displaystyle t_{0} \neq 0 $ basta con considerar
\[\psi\left(t\right)=\exp\left(A\left(t-t_{0}\right)\right) .\]
\end{observation}
\begin{prop}
\begin{enumerate}
\item Si $\displaystyle P\in \mathcal{M}_{n} $ es tal que $\displaystyle AP = PB $, entonces $\displaystyle e^{tA}P = Pe^{tB} $.
\item Si $\displaystyle AB = BA $, entonces $\displaystyle e^{tA}B = Be^{tA} $, $\displaystyle \forall t \in \R $. En particular, $\displaystyle AB = BA $. 
\item Si $\displaystyle AB = BA $, entonces $\displaystyle e^{t\left(A + B\right)}=e^{tA}e^{tB} $, $\displaystyle \forall t \in \R $. En particular, $\displaystyle e^{tA}e^{tB}=e^{tB}e^{tA} $, $\displaystyle \forall t \in \R $. 
\end{enumerate} 
\end{prop}
\begin{proof}
\begin{enumerate}
\item Podemos observar que 
	\[A^{2}P = A\left(AP\right) = A\left(PB\right)=\left(AP\right)B = PB^{2} .\]
A partir de aquí es sencillo demostrar por inducción que 
\[A^{k}P = PB^{k}, \; \forall k \in \N .\]
De esta forma, tenemos que
\[e^{tA}P = \sum^{\infty}_{k = 0}\frac{t ^{k}A^{k}}{k!}P = \sum^{\infty}_{k = 0}\frac{t ^{k}PB^{k}}{k!}=P\sum^{\infty}_{k = 0}\frac{t ^{k}B^{k}}{k!} = Pe^{tB} .\]
\item Basta tomar $\displaystyle B = A $ y $\displaystyle P = B $ y aplicar \textbf{(1)}. 
\item Si planteamos el PVI matricial
	\[
	\begin{cases}
	\phi'\left(t\right)=\left(A + B\right)\phi \\
	\phi\left(0\right)=I
	\end{cases}
	.\]
Tenemos que $\displaystyle \phi\left(t\right)=e^{t\left(A + B\right)} $ es la única solución (por el teorema de existencia y unicidad). Por otro lado, si tomamos $\displaystyle \psi\left(t\right)=e^{tA}e^{tB} $ tenemos que 
\[\psi'\left(t\right) = Ae^{tA}e^{tB}+e^{tA}Be^{tB} = Ae^{tA}e^{tB} + Be^{tA}e^{tB}=\left(A + B\right)e^{tA}e^{tB} = \left(A + B\right)\psi\left(t\right) .\]
Dado que $\displaystyle \psi\left(0\right)=I $, tenemos que por unicidad de soluciones debe ser que $\displaystyle \phi\left(t\right)=\psi\left(t\right) $. La última afirmación se deduce fácilmente a partir de la igualdad $\displaystyle e^{t\left(A + B\right)}=e^{t\left(B + A\right)} $. 
\end{enumerate}
\end{proof}
\begin{eg}
Consideremos el sistema $\displaystyle \vec{x}' = A \vec{x} $ con 
\[A = \begin{pmatrix} 0 & 1 & -1 \\ -3 & 4 & -1 \\ -1 & 1 & 2 \end{pmatrix} ,\]
y calculemos una matriz fundamental. Tendremos que el polinomio característico de $\displaystyle A $ es
\[P_{A}\left(\lambda \right)= \begin{vmatrix} -\lambda & 1 & -1 \\ -3 & 4-\lambda & -1 \\ -1 & 1 & 2-\lambda \end{vmatrix} =\left(\lambda-1\right)\left(\lambda-2\right)\left(\lambda-3\right).\]
Así, para cada autovalor se cumple que la multiplicidad geométrica es igual que la algebraica por lo que $\displaystyle A $ es diagonalizable. Tendremos que existe $\displaystyle P \in \mathcal{M}_{n}\left(\R\right) $ tal que 
\[A = P\begin{pmatrix} 1 & & \\ & 2 & \\ & & 3 \end{pmatrix}P^{-1} .\]
De esta forma, una matriz fundamental para $\displaystyle A $ será
\[\exp\left(At\right)=P\begin{pmatrix} e^{t} & & \\ & e^{2t} & \\ & & e^{3t} \end{pmatrix} .\]
\end{eg}

% --------------------------------------
\begin{observation}
Algunas observaciones antes de hacer la demostración.
\begin{enumerate}
\item \item Si $\displaystyle A $ no es diagonal podemos recurrir a la forma canónica de Jordan. Para calcular $\displaystyle \exp\left(At\right) $ me basta con calcular $\displaystyle \exp\left(B_{j}t\right) $, donde $\displaystyle B_{j} $ son los bloques de la matriz de Jordan de $\displaystyle A $. Podemos descomponer $\displaystyle B_{j}=D_{j}+N_{j} $, donde $\displaystyle N_{j} $ es la parte de debajo de la diagonal que es una matriz nilpontente. Nos gustaría que 
	\[\exp\left(B_{j}t\right)=\exp\left(t\left(D_{j}+N_{j}\right)\right)=?\exp\left(tD_{j}\right)\exp\left(tN_{j}\right) .\]
Una propiedad buena de la exponencial es que si $\displaystyle AB = BA $, entonces
\[\exp\left(t\left(A + B\right)\right)=\exp\left(tA\right)\exp\left(tB\right) .\]
Así, tendremos que 
\[e^{Bt}=\begin{pmatrix} e^{B_{1}t} & &0 & \\
	& e^{B_{2}t} & & \\
	& 0& & \\
	& & & e^{B_{n}t}\end{pmatrix} .\]
Así, el problema se reduce a calcular $\displaystyle e^{B_{i}t} $. Tendremos que 
\[e^{B_{i}t}=e^{\left(D_{i}+N_{i}\right)t}=e^{D_{i}t}e^{N_{i}t} .\]
De esta forma,
\[D_{i}=\begin{pmatrix} \lambda_{i} & &0 & \\
	& \lambda_{i} & & \\
	& 0& & \\
	& & & \lambda_{i}\end{pmatrix} \Rightarrow e^{D_{i}t} = \begin{pmatrix} e^{\lambda_{i}t} & &0 & \\
	& e^{\lambda_{i}t} & & \\
	& 0& & \\
	& & & e^{\lambda_{i}t}\end{pmatrix}.\]
Por otro lado, 
\[N_{i}=\begin{pmatrix} 0 & 1& & \\
	& 0 & 1& \\
	& & & 1\\
	& & & 0  \end{pmatrix} \Rightarrow e^{N_{i}t}=\sum^{\infty}_{k = 0}\frac{N_{i}^{k}t^{k}}{k !}= \sum^{l-1}_{k = 0}\frac{N_{i}^{k}t ^{k}}{k!}.\]
Es fácil ver que
	\[e^{N_{i}t}=\begin{pmatrix} 1 & t & \frac{t ^{2}}{2!} & \cdots & \frac{t ^{l-1}}{\left(l-1\right)!} \\
	& & & & \vdots \\
	& & & & \frac{t ^{2}}{2!} \\
	& & & & t \\
	& & & & 1\end{pmatrix} .\]

\item A parte de demostrar el teorema y las buenas propiedades de $\displaystyle \exp\left(At\right) $, falta ver qué pasa cuando $\displaystyle A $ tiene autovalores complejos. Supongamos que $\displaystyle \lambda = a + bi $ es un autovalor complejo. Entonces, tendremos que también es autovalor $\displaystyle \overline{\lambda }=a - bi $. El autovector asociado a $\displaystyle \lambda  $ será un vector de la forma $\displaystyle \vec{w}=\vec{u}+i\vec{v} $ y el vector asociado a $\displaystyle \overline{\lambda } $ será $\displaystyle \vec{\overline{w}}=\vec{u}-i\vec{v} $. Así, tendremos que on el razonamiento anterior $\displaystyle e^{\lambda t}\vec{w} $ y $\displaystyle e^{\overline{\lambda }t}\vec{\overline{w}} $ son soluciones pero están en $\displaystyle \C $. 
Entonces, tenemos que buscar una combinación lineal de ellas que sea real. 
\end{enumerate}
\end{observation}
\begin{prop}
\[\frac{d}{dt}\exp\left(At\right)=A \exp\left(At\right) .\]
\end{prop}
\begin{prop}
\begin{enumerate}
\item Si $\displaystyle P \in \mathcal{M}_{n}\left(\R\right) $ tal que $\displaystyle AP = PB $, entonces $\displaystyle e^{tA}P=Pe^{tB}  $.
\item Si $\displaystyle AB = BA $, entonces $\displaystyle e^{tA}B = Be^{tA} $. 
\item Si $\displaystyle AB = BA $, entonces $\displaystyle e^{t\left(A + B\right)}= e^{tA}e^{tB} $. 
\end{enumerate}
\end{prop}
\begin{proof}
Tenemos que (2) es trivial a partir de (1) si cogemos $\displaystyle P = B $ y $\displaystyle B = A $. Además, (1) es consecuencia de que
\[\exp\left(tA\right)=\sum^{\infty}_{k = 0}\frac{t ^{k}A^{k}}{k!} .\]
Así, tendremos que 
\[
\begin{split}
	e^{tA}P = & \sum^{\infty}_{k = 0}\frac{t ^{k}}{k!}A^{k}P = \sum^{\infty}_{k = 0}\frac{t ^{k}}{k!}PB^{k}=Pe^{tB}.
\end{split}
\]
Hemos usado que $\displaystyle A^{k}P = PB^{k} $. Para (3), puedo probarlo hacieno uso del teorema de existencia y unicidad. Podemos plantear el problema matricial
\[
\begin{cases}
\phi'=\left(A + B\right)\phi \\
\phi\left(0\right)=I
\end{cases}
.\]
Observamos que $\displaystyle \exp\left(t\left(A + B\right)\right) $ es solución y tiene que ser única, por el teorema de existencia y unicidad. Tenemos que $\displaystyle \varphi\left(t\right)=e^{tA}e^{tB} $, así
\[
\begin{cases}
\varphi'\left(t\right)= Ae^{tA} e^{tB} + e^{tA}Be^{tB} = Ae^{tA}e^{tB}+Be^{tA}e^{tB}\\
\varphi\left(0\right)=I
\end{cases}
.\]
\end{proof}
\section{Resumen}
Para resolver el problema
\[
\begin{cases}
\vec{x}'=A\vec{x} + \vec{b}\left(t\right) \\
\vec{x}\left(t_{0}\right)=\vec{x}_{0}
\end{cases}
.\]
Tendremos que la solución general será de la forma
\[\vec{x}\left(t\right)=\vec{x}_{h}\left(t\right)+\vec{x}_{p}\left(t\right) .\]
La solución particular la sacamos con el método de variación de constantes. Para calcular $\displaystyle \vec{x}_{h} $, tenemos que calcular $\displaystyle \psi $ para lo que necesitamos calcular el polinomio característico y calcular los autovalores y así sacar las cajas de Jordan. Si tenemos algún autovalor $\displaystyle \lambda \in \C $ con $\displaystyle \lambda = a + bi $, $\displaystyle a,b \in \R $. Tendremos que hay una solución para $\displaystyle \vec{x}' = A \vec{x} $ de la forma $\displaystyle \vec{z}\left(t\right)=e^{\lambda t}\vec{w} $, entonces $\displaystyle \vec{z} $ es una función compleja y nosotros necesitamos soluciones reales.
\begin{prop}
Consideramos un sistema liineal homogéneo $\displaystyle \vec{x}'=A\vec{x} $, con $\displaystyle A \in \mathcal{M}_{n}\left(\R\right) $, dada $\displaystyle z : \R \to \C^{n} $ tal que $\displaystyle \vec{z}\left(t\right)= \vec{x}\left(t\right)+\vec{y}\left(t\right)i $, con $\displaystyle \vec{x}, \vec{y} : \R \to \R^{n} $. Entonces, $\displaystyle \vec{z} $ es solución del sistema si y solo si $\displaystyle \vec{x} $ e $\displaystyle \vec{y} $ son soluciones del sistema. 
\end{prop}
Con esta proposición podemos hallar a partir de las soluciones complejas, soluciones reales y linealmente independientes. En efecto, si $\displaystyle \lambda  $ es un autovalor complejo de $\displaystyle A $, $\displaystyle \overline{\lambda } $ también lo será y si $\displaystyle \vec{w} $ es el autovector asociado a $\displaystyle \lambda  $, $\displaystyle \vec{\overline{w}} $ será el autovector asociado a $\displaystyle  \overline{\lambda } $. Así, tendremos que son soluciones
\[\vec{z}_{1}=e^{\lambda t}\vec{w} \quad \text{y} \quad \vec{z}_{2}=e^{\overline{\lambda }t}\vec{\overline{w}} .\]
Tenemos que existen $\displaystyle \vec{u}, \vec{v} $ tales que $\displaystyle \vec{w}= \vec{u}+\vec{v}i $. Así, tendremos que 
\[\vec{z}_{1}\left(t\right)=e^{\left(a+bi\right)t}\left(\vec{u}+i\vec{v}\right) = e^{at}\left(\cos bt + i \sin bt\right)\left(\vec{u}+i\vec{v}\right) .\]
\[\vec{z}_{2}\left(t\right)=e^{\left(a-bi\right)t}\left(\vec{u}-i\vec{v}\right)= e^{at}\left(\cos bt - i\sin bt\right)\left(\vec{u}-i\vec{v}\right) .\]
Cogemos 
\[\vec{x}\left(t\right)=\frac{\vec{z}_{1}\left(t\right)+\vec{z}_{2}\left(t\right)}{2} \quad \text{y} \quad \vec{y}\left(t\right)=\frac{\vec{z}_{1}\left(t\right)-\vec{z}_{2}\left(t\right)}{2} .\]
\section{El caso particular de las EDOS de orden $\displaystyle n $}
Consideremos una EDO lineal homogénea de orden $\displaystyle n $ con coeficientes constantes
\[x^{\left(n\right)} + a_{n-1}x^{\left(n-1\right)} + \cdots + a_{1}x' + a_{0}x = 0 .\]
Sabemos que podemos transformar esta EDO en un sistema 
\[
\begin{cases}
x_{1}=x \\
x_{2}=x' \\
\vdots \\
x_{n}=x^{\left(n-1\right)}
\end{cases}
\Rightarrow 
\begin{cases}
x_{1}' = x_{2} \\
x_{2}' = x_{3} \\
\vdots \\
x_{n-1}' = x_{n} \\
x_{n}' = -a_{n-1}x_{n} - \cdots - a_{1}x_{2}-a_{0}x_{1}
\end{cases}
.\]
Tendremos que 
\[A = \begin{pmatrix} 0 & 1 & 0 & \cdots & 0 \\
0 & 0 & 1 & \cdots & 0 \\
 \vdots & \vdots & \vdots & \vdots & \vdots  \\
0 & 0 & 0 & \cdots & 1 \\
-a_{0} & -a_{1} & - a_{2} & \cdots & - a_{n-1}\end{pmatrix} .\]
Para resolver $\displaystyle \left(S_{h}\right) $ calculamos $\displaystyle \phi\left(t\right)=e^{tA} $, para lo cual calculamos $\displaystyle P_{A}\left(\lambda \right)=\det\left(A-\lambda I\right) $. 
\begin{prop}
Para la matriz $\displaystyle A $ del sistema $\displaystyle \left(S_{h}\right) $ se cumple 
\[P_{A}\left(\lambda \right)=\left(-1\right)^{n}\left(\lambda^{n} + a_{n-1}\lambda^{n-1} + \cdots + a_{1}\lambda + a_{0}\right) .\]
Entonces, para calcular los autovalores bastará con calcular las raíces de lo que llamamos la \textbf{ecuación característica} 
\[ \left(E\right) \; \lambda^{n} + a_{n-1}\lambda^{n-1} + \cdots + a_{1}\lambda + a_{0}=0 .\]
\end{prop}
\begin{proof}
Se demuestra por inducción.
\end{proof}

\begin{observation}
Así, podemos asegurar con lo visto hasta ahora que $\displaystyle \lambda  $ es raíz de $\displaystyle \left(E\right) $ si y solo si $\displaystyle e^{\lambda t} $ es solución de la EDO. Nuestro objetivo es encontrar $\displaystyle n $ soluciones linealmente independintes, que formarán un conjunto fundamental de la EDO y la solución general será la combinación de las funciones (escalares) del conjunto fundamental. 
\end{observation}
\begin{prop}
La forma canónica de Jordan de $\displaystyle A $ dada en $\displaystyle \left(S_{h}\right) $, tiene un solo bloque elemental correspondiente a cada autovalor real o complejo de $\displaystyle A $. Por tanto, el conjunto fundamental estará formado por funciones del tipo 
\[ \left\{ e^{\lambda t}, t ^{k}e^{\lambda t}, e^{at}\cos bt, e^{at}\sin bt\right\}  .\]
\end{prop}
\begin{proof}
Queremos ver que $\displaystyle \forall \lambda  $ raíz de $\displaystyle P_{A}\left(\lambda \right)=0 $, $\displaystyle \dim\left(\Ker\left(A-\lambda I\right)\right)=1 $. Tenemos que 
\[ \left|A-\lambda I\right|=\begin{vmatrix} -\lambda & 1 & 0 & \cdots & 0\\
0 & -\lambda & 1 & \cdots & 0 \\
  & & & & \\
0 & 0 & 0 & \cdots & 1 \\
-a_{0} & -a_{1} & -a_{2} & \cdots & -a_{n-1}-\lambda\end{vmatrix}  .\]
Tenemos que $\displaystyle \ran\left(A-\lambda I\right)=n-1 $. Por tanto, 
\[\dim\left(\Ker\left(A -\lambda I\right)\right)= n-\ran\left(A - \lambda I\right)=n - \left(n-1\right)=1 .\]
\end{proof}

\begin{eg}
Calcula la solución general de la EDO
\[x^{\left(4\right)}-2x^{\left(3\right)} - 2x'' + 6x' + 5x = 0 .\]
Buscamos 4 soluciones linealmente independientes de la forma $\displaystyle e^{\lambda t} $, con $\displaystyle \lambda  $ raíz de
\[\lambda^{4} - 2\lambda^{3} - 2\lambda^{2}+6\lambda + 5 = 0 .\]
Las raíces son $\displaystyle \lambda \in \left\{ -1, 2 \pm i\right\}  $. Así, tenemos que el conjunto fundamental será $\displaystyle \left\{ e^{-t}, te^{-t}, e^{2t}\cos t, e^{2t}\sin t\right\}  $. La solución general es la familia cuatri-paramétrica
\[x\left(t\right)=C_{1}e^{-t} + C_{2}te^{-t} + C_{3}e^{2t}\cos t + C_{4}e^{2t}\sin t, \; C_{1}, C_{2}, C_{3}, C_{4} \in \R .\]
\end{eg}
\begin{eg}[Hoja 3, Ejercicio 1]
Dada 
\[\left(e^{Ax}\right)^{-1}=\begin{pmatrix} e^{-2x} & -xe^{-2x} \\ 0 & e^{-2x}\end{pmatrix} ,\]
con $\displaystyle A $ la matriz del sistema 
\[
\begin{cases}
y_{1}'= 2y_{1} + y_{2}+f\left(x\right) \\
y_{2}' = 2y_{2} + 4e^{2x}
\end{cases}
.\]
Además, una solución particular es 
\[\vec{y}\left(x\right)=\begin{pmatrix} \left(3x + 2x^{2}\right)e^{2x} \\ g\left(x\right) \end{pmatrix} .\]
Tenemos que calcular $\displaystyle f\left(x\right) $ y $\displaystyle g\left(x\right) $. Si no supiésemos nada sobre matrices fundamentales y exponenciales, podríamos calcular $\displaystyle g $ puesto que sabemos que 
\[g'\left(x\right)=2g\left(x\right) + 4e^{2x} .\]
Siguiendo el mismo razonamiento, habiendo calculado $\displaystyle g $ podemos sustituir los datos de la solución particular para calcular $\displaystyle f\left(x\right) $. \\
Intentemos resolverlo utilizando matrices fundamentales. Sabemos que la solución general del sistema será $\displaystyle \vec{y}\left(x\right)=\vec{y}_{h}\left(x\right) + \vec{y}_{p}\left(x\right) $, donde $\displaystyle \vec{y}_{h}\left(x\right)=e^{Ax}\vec{c} $ y $\displaystyle \vec{y}_{p}\left(x\right)=e^{Ax}\vec{c}\left(x\right) $. 
Tendremos que 
\[\left(e^{Ax}\vec{c}\left(x\right)\right)' = Ae^{Ax}\vec{c} + \begin{pmatrix} f\left(x\right) \\ 4e^{2x} \end{pmatrix} \iff e^{Ax}\vec{c}' = \begin{pmatrix} f\left(x\right) \\ 4e^{2x} \end{pmatrix}\iff \vec{c}'\left(x\right)=\left(e^{Ax}\right)^{-1}\begin{pmatrix} f\left(x\right) \\ 4e^{2x} \end{pmatrix} .\]
Tendremos que 
\[\vec{c}'\left(x\right)= \begin{pmatrix} e^{-2x} & -xe^{-2x} \\ 0 & e^{-2x}\end{pmatrix}\begin{pmatrix} f\left(x\right) \\ 4e^{2x} \end{pmatrix} = \begin{pmatrix} e^{-2x}f\left(x\right) -4x \\ 4 \end{pmatrix}.\]
Además, queremos que $\displaystyle e^{Ax}\vec{c} $ coincida con la solución particular que nos dan. Así, 
\[e^{Ax}\vec{c} = \begin{pmatrix} \left(3x + 2x^{2}\right)e^{2x} \\ g\left(x\right) \end{pmatrix} \iff \vec{c}=\left(e^{Ax}\right)^{-1}\begin{pmatrix} \left(3x + 2x^{2}\right)e^{2x} \\ g\left(x\right) \end{pmatrix} = \begin{pmatrix} 3x + 2x^{2} -xe^{-2x}g\left(x\right) \\ e^{-2x}g\left(x\right) \end{pmatrix} .\]
Si igualamos ambas expresiones tendremos que $\displaystyle 4x = e^{-2x}g\left(x\right) $, por lo que $\displaystyle g\left(x\right)=e^{2x}4x $ y tenemos que $\displaystyle e^{-2x}f\left(x\right)=3 $, por lo que $\displaystyle f\left(x\right)=3e^{2x} $. 
\end{eg}
\begin{eg}[Hoja 3, Ejercicio 2]
Demostrar que si $\displaystyle \phi $ es matriz fundamental principal para
\[\left(P\right)
\begin{cases}
\vec{x}' = A\left(t\right)\vec{x} \\
\vec{x}\left(t_{0}\right)=\vec{x}_{0}
\end{cases}
\]
entonces la solución general de $\displaystyle \left(P\right) $ es $\displaystyle \vec{x}\left(t\right)=\phi\left(t\right)\vec{x}_{0} $. Ya lo hemos visto en clase. Ahora, si $\displaystyle A\left(t\right)=A $, entonces la matriz fundamental principal tiene dos propiedades buenas:
\begin{enumerate}
\item $\displaystyle \phi\left(t + s\right)= \phi\left(t\right)\phi\left(s\right) $, $\displaystyle s,t \in \R $. 
\item $\displaystyle \phi^{-1}\left(t\right)=\phi\left(-t\right) $. 
\end{enumerate}
Claramente, (2) es consecuencia de (1) tomando $\displaystyle s = -t $. Para demostrar (1) utilizamos el teorema de existencia y unicidad. Tenemos que ver que tanto $\displaystyle \phi\left(t + s\right) $ como $\displaystyle \phi\left(t\right)\phi\left(s\right) $ son ambas soluciones de. problema de valor inicial
\[
\begin{cases}
\vec{X}' = A\vec{X} \\
\vec{X}\left(0\right)=\phi\left(s\right)
\end{cases}
,\]
donde $\displaystyle \vec{X} $ es una matriz. 
En efecto, tenemos que si derivamos con respecto a $\displaystyle t $ 
\[\left(\phi\left(t\right)\phi\left(s\right)\right)' = \phi'\left(t\right)\phi\left(s\right) = A\phi\left(t\right)\phi\left(s\right).\]
Así, es solución y además $\displaystyle \phi\left(0\right)\phi\left(s\right)=\phi\left(s\right) $, por lo que es solución del problema. También $\displaystyle \phi\left(t + s\right) $ verifica la condición inicial: $\displaystyle \phi\left(0 + s\right)= \phi\left(s\right) $. Además, 
\[\phi'\left(t + s\right)=A\phi\left(t + s\right) .\]
Por el teorema de existencia y unicidad debe ser que $\displaystyle \phi\left(t + s\right)=\phi\left(t\right)\phi\left(s\right) $. Merece la pena comprobar por qué es importante la condición de que $\displaystyle A $ sea de coeficientes constantes. 
\end{eg}
\begin{eg}[Hoja 3, Ejercicio 3]
Sea
\[\phi\left(t\right)=\begin{pmatrix} 0 & e^{3t} & \frac{1}{t+1} \\ t+1 & 0 & e^{-3t} \\ 1 & 0 & 0 \end{pmatrix} .\]
Discutir $\displaystyle t $ tal que $\displaystyle \phi\left(t\right) $ es solución de algún sistema y calcular la solución general de este. Para $\displaystyle t \neq -1 $ tiene snetido el problema. Como el determinante es no nulo $\displaystyle \forall t \in \left(-1, \infty\right) $, tenemos que puede ser matriz fundamental. 
\end{eg}
\section{Sistemas planos y diagramas de fase}
Vamos a estudiar cualitativamente los sistemas $\displaystyle \vec{x}' = A\vec{x} $. 
